\chapter{How a MERN Application Works}

Every feature in this book is a variation of the same round trip: a user
action in the browser eventually changes data in the database and the UI
updates to reflect it.

\section{The Full Request Lifecycle}

A MERN request always travels through the same layers, in the same order,
whether it's a login form or a task creation button.

\begin{diagrambox}[MERN Request Lifecycle]
\begin{verbatim}
 User -> React Component (JSX + handler)
      -> Event Handler (onClick/onSubmit)
      -> API function (services/taskService.js)
      -> Axios instance (src/api/axios.js)
      -> HTTP Request (method + URL + headers + body)
      -> Express App -> Router (routes/taskRoutes.js)
      -> Middleware (auth, validation)
      -> Controller (controllers/taskController.js)
      -> Service/Model layer
      -> Mongoose (Task.find / Task.create / ...)
      -> MongoDB (data on disk)
      -> back up through Controller ({ success, data })
      -> JSON Response
      -> Axios (resolves/rejects)
      -> React State (useState/context)
      -> UI Update (re-render)
\end{verbatim}
\end{diagrambox}

\begin{itemize}
  \item \textbf{React Component} -- renders UI, listens for interaction.
  \item \textbf{Event Handler} -- captures interaction, decides what to call.
  \item \textbf{API function} -- wrapper knowing the endpoint and payload shape.
  \item \textbf{Axios} -- HTTP client; attaches base URL, headers, credentials.
  \item \textbf{Express App/Router} -- matches the request to a route.
  \item \textbf{Middleware} -- cross-cutting checks before the controller.
  \item \textbf{Controller} -- validates input, calls the model, builds response.
  \item \textbf{Mongoose} -- translates JS calls into MongoDB queries and back.
  \item \textbf{MongoDB} -- persists and retrieves documents.
  \item \textbf{JSON Response} -- consistent \texttt{\{ success, data \}} envelope.
  \item \textbf{React State} -- single source of truth driving the UI.
\end{itemize}

\begin{notebox}[NOTE]
Every chapter teaches one segment of this diagram in depth -- the chain
itself never changes as the app grows.
\end{notebox}

\section{Concrete Example: ``User Clicks Create Task''}

Tracing one real interaction end to end in the \textbf{Task Management System}:

\begin{diagrambox}[Example: Create Task]
\begin{verbatim}
1. User fills Title+Description, clicks "Create Task"
2. TaskForm.jsx -> handleSubmit(e)
3. taskService.js -> createTask({ title, description })
4. axios -> POST /api/tasks (JSON body, withCredentials)
5. Express router -> router.post("/tasks", protect, createTask)
6. protect middleware -> verifies JWT cookie, sets req.user
7. createTask controller -> validates body, Task.create({...})
8. Mongoose inserts document into "tasks" collection
9. MongoDB returns new document with _id
10. Controller -> res.status(201).json({ success: true, data: task })
11. Axios promise resolves with response.data
12. React setTasks -> new task appended to state
13. UI -> TaskList re-renders, new TaskCard appears
\end{verbatim}
\end{diagrambox}

\begin{enumerate}
  \item \textbf{User action:} types into two inputs, clicks submit -- no network yet.
  \item \textbf{Component:} \texttt{TaskForm.jsx} owns form state via \texttt{useState}.
  \item \textbf{Service function:} \texttt{taskService.js} is the only place that knows the endpoint URL.
  \item \textbf{Axios:} shared instance adds \texttt{baseURL}, \texttt{Content-Type}, cookies.
  \item \textbf{Express router:} \texttt{/api/tasks} maps to \texttt{createTask}.
  \item \textbf{Middleware:} \texttt{protect} rejects with 401 before any DB access.
  \item \textbf{Controller:} validates payload, delegates persistence to Mongoose.
  \item \textbf{Mongoose/MongoDB:} writes the document, returns it with \texttt{\_id}.
  \item \textbf{Response:} controller always replies \texttt{\{ success, data \}} (Chapter 8).
  \item \textbf{Frontend update:} resolved data is pushed into state; list re-renders.
\end{enumerate}

\begin{tipbox}[TIP]
When debugging, walk this chain link by link -- most bugs are a mismatch
between two adjacent links (wrong URL, wrong field name, missing
\texttt{await}, wrong HTTP method).
\end{tipbox}
